\section{Divisor exponents of genuine ideal powers}
\label{sec:hironaka-divisor-extension}

The monomial factor in W{\l}odarczyk's Step~2 is a product of principal
divisor ideals, with multiplicities.  The following local algebra records
how those multiplicities change when an ideal is replaced by a genuine
ideal power.  It is an extension of the order and divisor-exponent calculus,
not a construction of the global decomposition of every principalizing
divisor into the members of a prescribed simple normal crossings family.

\begin{lemma}[Successor thresholds for ideal powers]
\label{hir:divisor-power-successor-threshold}
\uses{hir:order-power,hir:order-basic,hir:cosupport}
\lean{Hironaka.HironakaDivisorPowers.cosupport_idealPow_mul_succ_eq}
\leanok
Let $\mathfrak a$ be an ideal of holomorphic germs on a complex manifold
$M$, and let $k\geq 1$ and $e\geq 0$ be integers.  Then
\[
 \operatorname{cosupp}(\mathfrak a^k,ke+1)
 =\operatorname{cosupp}(\mathfrak a,e+1).
\]
\end{lemma}
\begin{proof}
If every germ of $\mathfrak a_x$ has order at least $e+1$, then every
product of $k$ such germs has order at least $k(e+1)$.  Finite sums and
holomorphic multiples preserve this lower bound.  Since
$ke+1\leq k(e+1)$, every germ of $(\mathfrak a^k)_x$ has order at least
$ke+1$.  Conversely, if a germ $g\in\mathfrak a_x$ has order at most $e$,
then multiplicativity of the order gives
$\operatorname{ord}_x(g^k)\leq ke$.  The germ $g^k$ belongs to
$(\mathfrak a^k)_x$, so this stalk does not have order at least $ke+1$.
These two implications prove the equality.
\end{proof}

\begin{lemma}[Divisibility thresholds for ideal powers]
\label{hir:divisor-power-divisibility-threshold}
\uses{hir:divisor-power-successor-threshold,hir:divisor-exponent}
\lean{Hironaka.HironakaDivisorPowers.isDivNear_idealPow_mul_succ_iff}
\leanok
For a subset $H\subset M$, a point $x\in M$, $k\geq1$ and $e\geq0$,
the ideal $\mathfrak a^k$ has order at least $ke+1$ along $H$ near $x$
if and only if $\mathfrak a$ has order at least $e+1$ along $H$ near $x$.
\end{lemma}
\begin{proof}
Apply Lemma~\ref{hir:divisor-power-successor-threshold} at each point
of $H$ in a neighbourhood of $x$.  The equality of the two cosupports
identifies the two neighbourhood conditions in both directions.
\end{proof}

\begin{lemma}[A power of a nonzero ideal stalk is nonzero]
\label{hir:divisor-power-nonzero}
\uses{hir:order-power,hir:divisor-exponent}
\lean{Hironaka.HironakaDivisorPowers.isNonzeroAt_idealPow}
\leanok
If $\mathfrak a_x$ contains a holomorphic germ which does not vanish
identically near $x$, then $(\mathfrak a^k)_x$ has the same property for
every integer $k\geq0$.
\end{lemma}
\begin{proof}
Choose such a germ $g$.  Its power $g^k$ lies in $(\mathfrak a^k)_x$.
At each point where $g$ is nonzero, $g^k$ is nonzero; these points occur
arbitrarily close to $x$.  This also covers $k=0$, since $g^0=1$.
\end{proof}

\begin{theorem}[Divisor exponents scale with ideal powers]
\label{hir:divisor-exponent-ideal-power}
\uses{hir:divisor-power-divisibility-threshold,hir:divisor-power-nonzero,hir:divisor-exponent}
\dcref{W{\l}odarczyk 2005, Step 2}
\source{wlodarczyk05-hironaka}
\lean{Hironaka.HironakaDivisorPowers.divExp_idealPow}
\leanok
Let $x\in H\subset M$, and let $\mathfrak a$ be an ideal of holomorphic
germs whose stalk at $x$ is nonzero.  For every integer $k\geq0$,
\[
 e_H(\mathfrak a^k,x)=k\,e_H(\mathfrak a,x).
\]
\end{theorem}
\begin{proof}
Put $e=e_H(\mathfrak a,x)$.  Nonvanishing makes this exponent finite:
$\mathfrak a$ has order at least $e$ along $H$ near $x$, and it does not
have order at least $e+1$ there.  Products of germs give divisibility of
$\mathfrak a^k$ to order $ke$, so its exponent is at least $ke$.
For $k\geq1$, Lemma~\ref{hir:divisor-power-divisibility-threshold}
rules out divisibility to order $ke+1$.  Lemma~\ref{hir:divisor-power-nonzero}
ensures that its exponent is also finite, hence is at most $ke$.
For $k=0$, the ideal is the unit ideal and its germ $1$ has order zero
at every point of $H$; its exponent is therefore zero.
\end{proof}

\begin{lemma}[Successor thresholds along a submanifold]
\label{hir:divisor-power-successor-threshold-along}
\uses{hir:divisor-power-successor-threshold,hir:exponent-along,hir:order-basic}
\lean{Hironaka.HironakaDivisorPowers.cosupport_idealPow_mul_succ_eq_along,Hironaka.HironakaDivisorPowers.isDivNearAlong_idealPow_mul_succ_iff}
\leanok
Let $N$ be a submanifold of $M$. For $k\geq1$ and $e\geq0$,
\[
 \operatorname{cosupp}(\mathfrak a^k|_N,ke+1)
 =\operatorname{cosupp}(\mathfrak a|_N,e+1).
\]
Consequently the corresponding neighbourhood divisibility conditions
along $N\cap H$ agree.
\end{lemma}
\begin{proof}
The inclusion from right to left follows by multiplying germs vanishing
to order $e+1$ on $N$, and using $ke+1\leq k(e+1)$.
For the reverse inclusion, fix $x\in N$ and a holomorphic retraction
$r$ onto $N$ near $x$.  A germ $g$ has order at least $q$ on $N$ at
$x$ if and only if $g\circ r$ has ambient order at least $q$.
If $g$ fails to have restricted order $e+1$, then
$g\circ r$ has ambient order at most $e$, and its $k$-th power has
ambient order at most $ke$.  Since
$(g^k)\circ r=(g\circ r)^k$, the germ $g^k\in\mathfrak a^k_x$ fails
to have restricted order $ke+1$.  This proves the reverse inclusion.
Applying the equality pointwise near $x$ proves the divisibility assertion.
\end{proof}

\begin{lemma}[Restricted nonvanishing survives powers]
\label{hir:divisor-power-nonzero-along}
\uses{hir:divisor-power-nonzero,hir:exponent-along}
\lean{Hironaka.HironakaDivisorPowers.isNonzeroAlongAt_idealPow}
\leanok
If $\mathfrak a_x$ contains a holomorphic germ whose restriction does
not vanish identically on $N$ near $x$, then $(\mathfrak a^k)_x$ has
the same property, for every $k\geq0$.
\end{lemma}
\begin{proof}
Choose the germ $g$ in the assumption.  Then $g^k$ belongs to
$(\mathfrak a^k)_x$ and is nonzero at each point of $N$ where $g$ is
nonzero. Such points occur arbitrarily close to $x$.
\end{proof}

\begin{theorem}[Divisor exponents of powers along a submanifold]
\label{hir:divisor-exponent-ideal-power-along}
\uses{hir:divisor-power-successor-threshold-along,hir:divisor-power-nonzero-along,hir:exponent-along}
\dcref{W{\l}odarczyk 2005, Step 2}
\source{wlodarczyk05-hironaka}
\lean{Hironaka.HironakaDivisorPowers.divExpAlong_idealPow}
\leanok
Let $N$ be a submanifold, $x\in N\cap H$, and suppose
$\mathfrak a_x|_N$ is nonzero. For every $k\geq0$,
\[
 e_{H|_N}(\mathfrak a^k,x)=k\,e_{H|_N}(\mathfrak a,x).
\]
\end{theorem}
\begin{proof}
Let $e=e_{H|_N}(\mathfrak a,x)$.  A local holomorphic retraction
onto $N$ and nonvanishing ensure that this exponent is finite.
Products give restricted divisibility of $\mathfrak a^k$ to order $ke$.
For $k\geq1$, Lemma~\ref{hir:divisor-power-successor-threshold-along}
excludes restricted divisibility to order $ke+1$, since it would imply
divisibility of $\mathfrak a$ to order $e+1$. Restricted nonvanishing
survives powers, so the exponent of $\mathfrak a^k$ is finite and equals
$ke$. When $k=0$, its unit germ has order zero on $N$, which gives
the asserted equality.
\end{proof}

\begin{lemma}[An invertible divisor ideal has nonzero stalks]
\label{hir:divisor-ideal-nonzero}
\uses{ana:ideal-sheaves,hir:divisor-exponent}
\lean{Hironaka.HironakaDivisorPowers.isNonzeroAt_of_isInvertible}
\leanok
Every stalk of an invertible divisor ideal contains a holomorphic germ
which does not vanish identically near that point.
\end{lemma}
\begin{proof}
Near a point $x$, choose a principal generator $g$ for the ideal.
Invertibility includes the fact that $g$ vanishes identically on no
nonempty open subset of that neighbourhood. Every open neighbourhood
of $x$ therefore contains a point where $g$ is nonzero. The germ $g$
belongs to the ideal stalk at $x$ and has the asserted property.
\end{proof}

\begin{corollary}[Multiplicities of powers of an actual divisor ideal]
\label{hir:divisor-ideal-power-multiplicity}
\uses{hir:divisor-ideal-nonzero,hir:divisor-exponent-ideal-power}
\lean{Hironaka.HironakaDivisorPowers.divExp_divisor_idealPow}
\leanok
For an effective divisor $D$ on $M$, $x\in H$, and $k\geq0$,
\[
 e_H(\mathcal O(-D)^k,x)=k\,e_H(\mathcal O(-D),x).
\]
\end{corollary}
\begin{proof}
Apply Theorem~\ref{hir:divisor-exponent-ideal-power} to the actual
invertible ideal $\mathcal O(-D)$; its holomorphic nonzero stalk members
are supplied by Lemma~\ref{hir:divisor-ideal-nonzero}.
\end{proof}

\subsection{Removing a genuine connected smooth component}

The next results construct one global multiplicity and a residual
effective divisor from the actual invertible ideal. The component is
assumed connected: a possibly disconnected member of a simple normal
crossings list need not carry one global multiplicity.

\begin{lemma}[Local constancy of an ideal exponent on a smooth carrier]
\label{hir:divisor-exponent-locally-constant-carrier}
\uses{hir:divisor-exponent}
\lean{Hironaka.HironakaDivisorResidual.isLocallyConstant_divExp_on_submanifold}
\leanok
Let $W\subset M$ be a closed submanifold and $\mathfrak a$ a coherent
ideal nonzero at every stalk. The function
$x\mapsto e_W(\mathfrak a,x)$ on $W$ is locally constant.
\end{lemma}
\begin{proof}
At each $x\in W$, the regularity of $W$ and nonvanishing of
$\mathfrak a$ in a neighbourhood give local constancy of the divisor
exponent on $W$. Restrict this neighbourhood to the subspace $W$.
Thus every point of $W$ has a neighbourhood on which the displayed
function equals its value at that point.
\end{proof}

\begin{lemma}[A connected carrier supplies one actual multiplicity]
\label{hir:divisor-exponent-connected-carrier}
\uses{hir:divisor-exponent-locally-constant-carrier,hir:divisor-exponent}
\lean{Hironaka.HironakaDivisorResidual.exists_divExp_on_connected_submanifold}
\leanok
Let $W$ be a nonempty connected closed submanifold of fixed dimension
and $\mathfrak a$ a coherent ideal nonzero at every stalk. There is
$a\in\mathbb N$ such that, for every $x\in W$ and $r\geq0$,
\[
 \mathfrak a\text{ has order at least }r\text{ along }W
 \text{ near }x \quad\Longleftrightarrow\quad r\leq a.
\]
\end{lemma}
\begin{proof}
A locally constant function on a connected space is constant: each
nonempty fibre is both open and closed, so its complement is empty.
Choose $p\in W$ and put $a=e_W(\mathfrak a,p)$. By
Lemma~\ref{hir:divisor-exponent-locally-constant-carrier},
$e_W(\mathfrak a,x)=a$ for every $x\in W$. Nonzero stalks ensure
that the finite exponent characterizes divisibility, proving the
displayed equivalence.
\end{proof}

\begin{lemma}[A smooth hypersurface has empty interior]
\label{hir:divisor-smooth-carrier-empty-interior}
\uses{hir:snc}
\lean{Hironaka.HironakaDivisorResidual.interior_eq_empty_of_smooth_hypersurface}
\leanok
If $W$ is a closed smooth submanifold of dimension $d$ in an
$n$-dimensional complex manifold and $d+1=n$, then $W$ has empty
interior.
\end{lemma}
\begin{proof}
An interior point of $W$ has an open neighbourhood contained in $W$;
in the ambient chart it is therefore regular of dimension $n$.
But every point of $W$ is regular of dimension $d$. The dimension
of a regular germ is unique, so $d=n$, contrary to $d+1=n$.
\end{proof}

\begin{theorem}[Actual residual after removing a connected smooth component]
\label{hir:divisor-connected-component-residual}
\uses{hir:divisor-exponent-connected-carrier,hir:divisor-smooth-carrier-empty-interior,hir:divisor-ideal-nonzero,hir:hadamard-bridge}
\dcref{W{\l}odarczyk 2005, Step 2}
\source{wlodarczyk05-hironaka}
\lean{Hironaka.HironakaDivisorResidual.exists_residual_on_connected_hypersurface}
\leanok
Let $D$ be an effective divisor on an $n$-dimensional complex
manifold $M$, and let $W$ be a nonempty connected closed smooth
submanifold of dimension $d$ with $d+1=n$. There exist $a\geq0$
and an effective divisor $R$ on $M$ such that
\[
 \begin{aligned}
 e_W(\mathcal O(-D),x)&=a &&(x\in W),\\
 \mathcal O(-D)&=\mathcal I_W^a\mathcal O(-R),\\
 \mathcal O(-R)&=(\mathcal O(-D):\mathcal I_W^a),
 \end{aligned}
\]
and $W$ is not contained in the support of $R$.
\end{theorem}
\begin{proof}
The ideal $\mathcal O(-D)$ is coherent and has nonzero stalks.
Lemma~\ref{hir:divisor-exponent-connected-carrier} therefore supplies
the constant finite exponent $a$. At every $x\in W$, all germs of
$\mathcal O(-D)_x$ vanish to order at least $a$ on $W$. The Hadamard
bridge gives the global inclusion
$\mathcal O(-D)\subset\mathcal I_W^a$, including the stalks off $W$.

The smooth hypersurface ideal $\mathcal I_W$ is invertible: locally
on $W$ it is generated by one flattening coordinate, and off $W$
by $1$; its local generator vanishes identically on no open subset
because $W$ has empty interior. Its power $\mathcal I_W^a$ is
invertible as well. The colon ideal
$\mathfrak r=(\mathcal O(-D):\mathcal I_W^a)$ is invertible.
Indeed, for local generators $u$ and $v$ of these two nested
invertible ideals, the inclusion gives $u=qv$ with $q$ holomorphic.
The germ $q$ generates the colon ideal: cancellation by $v$ proves
one inclusion, and multiplication by $v$ proves the other.
Nonvanishing of $u$ on every nonempty open set forces the same
nonvanishing for $q$. Thus $\mathfrak r=\mathcal O(-R)$ for an
effective divisor $R$, and the same cancellation yields
$\mathcal O(-D)=\mathcal I_W^a\mathcal O(-R)$.

If $W$ were contained in the support of $R$, every germ of the
coherent ideal $\mathcal O(-R)$ would vanish on $W$ near its base
point: membership propagates to nearby stalks, and each such germ
vanishes at the points of $W$. Hence
$\mathcal O(-R)\subset\mathcal I_W$, and the product identity would
give $\mathcal O(-D)\subset\mathcal I_W^{a+1}$.
At any $p\in W$ this implies divisibility to order $a+1$, contradicting
the characterization of the exponent $a$. Therefore $W$ is not
contained in the support of $R$.
\end{proof}

\begin{corollary}[Powers of the actual component factorization]
\label{hir:divisor-connected-component-residual-power}
\uses{hir:divisor-connected-component-residual}
\lean{Hironaka.HironakaDivisorResidual.exists_residual_power_on_connected_hypersurface}
\leanok
Under the hypotheses of Theorem~\ref{hir:divisor-connected-component-residual},
for each $k\geq0$ the same actual $a$ and residual $R$ satisfy
\[
 \mathcal O(-D)^k=\mathcal I_W^{ak}\mathcal O(-R)^k.
\]
\end{corollary}
\begin{proof}
Construct $a$ and $R$ by the preceding theorem, then take the $k$-th
power of its whole ideal identity. Products of coherent ideals commute
and associate, so the $k$-th power of a product is the product of the
$k$-th powers, and $(\mathcal I_W^a)^k=\mathcal I_W^{ak}$.
The exponent zero gives the unit ideal on both sides.
\end{proof}

\begin{definition}[Finite product of component powers]
\label{hir:divisor-finite-component-product}
\uses{ana:ideal-sheaves}
\lean{Hironaka.HironakaDivisorResidual.componentPowerProduct}
\leanok
For a finite ordered list $A=((W_1,a_1),\ldots,(W_s,a_s))$ of subsets
of $M$ and nonnegative integers, put
\[
 P(A)=\mathcal I_{W_1}^{a_1}\cdots\mathcal I_{W_s}^{a_s},\qquad
 P(\varnothing)=\mathcal O_M.
\]
\end{definition}

\begin{lemma}[Coherence of a finite component product]
\label{hir:divisor-finite-component-product-coherent}
\uses{hir:divisor-finite-component-product}
\lean{Hironaka.HironakaDivisorResidual.isCoherent_componentPowerProduct}
\leanok
If every $W_i$ is a closed submanifold, then $P(A)$ is coherent.
\end{lemma}
\begin{proof}
The vanishing ideal of each closed submanifold is coherent. Powers
and products of coherent ideals are coherent. Induction on the list,
starting from the coherent unit ideal, proves the claim.
\end{proof}

\begin{theorem}[Extraction along a finite family of connected hypersurfaces]
\label{hir:divisor-finite-connected-component-residual}
\uses{hir:divisor-connected-component-residual,hir:divisor-finite-component-product-coherent}
\lean{Hironaka.HironakaDivisorResidual.exists_residual_on_finite_connected_hypersurfaces}
\leanok
Let $E=(W_1,\ldots,W_s)$ be a finite list of nonempty connected closed
smooth hypersurfaces in $M$, and $D$ an effective divisor. There exist
nonnegative integers $a_1,\ldots,a_s$ and an effective divisor $R$ such that
\[
 \mathcal O(-D)=\mathcal I_{W_1}^{a_1}\cdots
 \mathcal I_{W_s}^{a_s}\mathcal O(-R),
 \qquad \operatorname{Supp}(R)\subset\operatorname{Supp}(D),
\]
and no $W_i$ is contained in $\operatorname{Supp}(R)$.
The list may be empty and may contain repeated members.
\end{theorem}
\begin{proof}
Induct on the length. For the empty list take $R=D$; the product is
the unit ideal. For $E=(W,E')$, Theorem~\ref{hir:divisor-connected-component-residual}
produces $a$ and $D_1$ with
$\mathcal O(-D)=\mathcal I_W^a\mathcal O(-D_1)$ and
$W\not\subset\operatorname{Supp}(D_1)$. The induction hypothesis
applied to $D_1$ and $E'$ produces $A'$ and $R$ with
$\mathcal O(-D_1)=P(A')\mathcal O(-R)$,
$\operatorname{Supp}(R)\subset\operatorname{Supp}(D_1)$, and no
member of $E'$ contained in $\operatorname{Supp}(R)$.
Associativity of coherent ideal products gives
$\mathcal O(-D)=(\mathcal I_W^aP(A'))\mathcal O(-R)$.
A zero of every germ of $\mathcal O(-D_1)$ is a zero of every germ
of its product with $\mathcal I_W^a$, so
$\operatorname{Supp}(D_1)\subset\operatorname{Supp}(D)$.
These inclusions give the asserted support inclusion and also rule
out $W\subset\operatorname{Supp}(R)$.
The members of $E'$ are handled by the induction hypothesis. The
argument does not require the members to be distinct, so repeated
members satisfy the same conclusion.
\end{proof}

\begin{remark}[Scope of the multiplicity calculation]
Each exponent extracted from a single connected component is defined
by neighbourhood divisibility of stalks. In the finite theorem, $a_i$
is the exponent of the intermediate divisor after the preceding members
have been removed; these exponents follow the order of the list. They
are not asserted to be the exponents of the original $D$ for every list
position. In particular a repeated member has exponent zero after its
first removal. For a smooth hypersurface the local exponent agrees with the largest
power of its vanishing ideal dividing the stalk.  Constructing a global
identity $\mathcal O(-D)=\prod_i\mathcal I_{E_i}^{a_i}$ for a general
simple-normal-crossings-supported divisor $D$, and then identifying
$D=\sum_i a_iE_i$, remains an additional obligation.
The finite extraction theorem leaves an actual residual factor. It
does not prove this factor is the unit ideal from an SNC support
assumption, nor does it refine disconnected SNC members into their
connected components.
\end{remark}
